
\paragraph{\textbf{Proposition 2.9.}}
\SGARefTarget{sga3:xiv:2.9}
{\itshape Let \(S,\mathfrak g\) be as in \SGARefLink{sga3:xiv:2.4}{2.4}. Consider the following
conditions:

\smallskip
\noindent
\textup{\((\mathrm C_0)\)}
The nilpotent rank of \(\mathfrak g(s)\), for \(s\in S\), is a
locally constant function of \(s\).

\smallskip
\noindent
\textup{\((\mathrm C_1)\)}
Locally for the fpqc topology, there exists a Cartan subalgebra of
\(\mathfrak g\).

\smallskip
\noindent
\textup{\((\mathrm C'_1)\)}
The same as \((\mathrm C_1)\), with ``fpqc topology'' replaced by
``étale topology''.

\smallskip
\noindent
\textup{\((\mathrm C_2)\)}
Condition \((\mathrm C_0)\) holds, and for every \(S'\) over \(S\),
every quasi-regular section of
\[
  \mathfrak g_{S'}=\mathfrak g\otimes_S S'
\]
is regular.

\smallskip
\noindent
\textup{\((\mathrm C_3)\)}
Every \(s\in S\) has an open neighborhood \(V\) on which the
Killing polynomial of \(\mathfrak g\) has the form
\[
  t^r\bigl(t^{n-r}+c_1t^{n-r-1}+\cdots+c_{n-r}\bigr),
\]
where, for every \(s\in V\),
\[
  c_{n-r}(s)\in
  \operatorname{Sym}\!\bigl(\mathfrak g(s)^\vee\bigr)
\]
is nonzero.

\smallskip
With this notation, one has the implications
\[
  (\mathrm C_3)\Longrightarrow(\mathrm C_2)
  \Longrightarrow(\mathrm C'_1)\Longrightarrow(\mathrm C_1)
  \Longrightarrow(\mathrm C_0),
\]
and, when \(S\) is reduced, the five conditions above are
equivalent.}

Notice also that these conditions are manifestly stable under base
change and are local for the fpqc topology.

The implications
\[
  (\mathrm C'_1)\Longrightarrow(\mathrm C_1)
  \Longrightarrow(\mathrm C_0)
\]
are trivial, and the implication
\((\mathrm C_0)\Longrightarrow(\mathrm C_3)\) when \(S\) is
reduced is immediate. We note moreover:

\paragraph{\textbf{Corollary 2.10.}}
\SGARefTarget{sga3:xiv:2.10}
\emph{Suppose that condition \((\mathrm C_0)\) holds. Let \(U\) be
the set of elements of \(W(\mathfrak g)\) that are regular in their
fiber. Then \(U\) is open. In particular, for every section \(a\)
of \(\mathfrak g\) over \(S\), the set \(V\) of \(s\in S\) for
which \(a(s)\in\mathfrak g(s)\) is regular is open.}

Indeed, the first assertion follows from the second, applied to
\(\mathfrak g_{S'}\) for every base change \(S'\) over \(S\). For
the second, we may suppose that \(S\) is reduced; then
\((\mathrm C_3)\) holds, and it is enough to inspect the Killing
polynomial of \(a\) in \(\mathfrak g\).

The implication
\((\mathrm C_2)\Longrightarrow(\mathrm C'_1)\) now follows easily
from part~b) of the following more precise statement.

\paragraph{\textbf{Corollary 2.11.}}
\SGARefTarget{sga3:xiv:2.11}
{\itshape Suppose that condition \((\mathrm C_2)\) holds. Then:

\smallskip
\noindent
\textup{a)}
For every \(s\in S\) and every Cartan subalgebra
\(\mathfrak d_0\) of \(\mathfrak g(s)\) such that
\(\mathfrak d_0\) contains a regular element of \(\mathfrak g(s)\)
\textup{(a condition automatically satisfied when \(k(s)\) is
infinite)}, there exist an open neighborhood \(V\) of \(s\) and a
Cartan subalgebra \(\mathfrak d\) of \(\mathfrak g|_V\) whose fiber
at \(s\) is \(\mathfrak d_0\). If \(S_1\) is a subscheme of \(S\)
containing \(s\), and if \(\mathfrak d_0\) has already been
extended to a Cartan subalgebra \(\mathfrak d_1\) of
\(\mathfrak g\otimes_S S_1\), then one can find an open
neighborhood \(V\) of \(s\) in \(S\) and a Cartan subalgebra
\(\mathfrak d\) of \(\mathfrak g|_V\) such that
\[
  \mathfrak d\otimes_V(S_1\cap V)
  =\mathfrak d_1|_{S_1\cap V}.
\]

\smallskip
\noindent
\textup{b)}
For every \(s\in S\) such that \(\mathfrak g(s)\) contains a
regular element \textup{(a condition automatically satisfied when
\(k(s)\) is infinite)}, there exist an open neighborhood \(V\) of
\(s\) and a Cartan subalgebra \(\mathfrak d\) of
\(\mathfrak g|_V\).}

Part~b) follows from part~a) by taking
\[
  \mathfrak d_0=\Nil\!\left(a_0,\mathfrak g(s)\right),
\]
where \(a_0\) is a regular element of \(\mathfrak g(s)\).
To prove part~a), in its second formulation, choose a regular
element \(a_0\) of \(\mathfrak g(s)\) contained in
\(\mathfrak d_0\). Extend it near \(s\) in \(S_1\) to a section of
\(\mathfrak d_1\), and extend that section in turn near \(s\) to a
section of \(\mathfrak g\). By \SGARefLink{sga3:xiv:2.10}{2.10} this section is quasi-regular
on an open neighborhood \(V\) of \(s\), hence regular by
\((\mathrm C_2)\). Therefore, by \SGARefLink{sga3:xiv:2.6}{2.6},
\[
  \mathfrak d=\Nil\!\left(a,\mathfrak g|_V\right)
\]
has the required properties.
